% TOP_PROBLEM: {"id": 4600050, "title": "Classify sofic shifts", "category_id": 11, "category": {"id": 11, "name": "dynamical_systems", "display_name": "Dynamical Systems", "description": "Problems about long-term behavior of deterministic systems, Hamiltonian dynamics, and periodic orbits.", "slug": "dynamical-systems", "order_index": 11, "created_at": "2026-07-31T15:26:25.671Z"}, "status": "partially_solved", "problem_number": "AMR-045-0050", "set_id": 13}
% FIRST_POSTED: 2026-10-01
% ENTRY_KIND: statement_only
% UnsolvedMath ID 4600050; source catalogue code AMR-045-0050
% !TeX program = lualatex
% Definition and sources only; this file is not an LLM solution attempt.
\documentclass[11pt,a4paper]{article}
\usepackage[margin=25mm]{geometry}
\usepackage{fontspec}
\setmainfont{lmroman10-regular.otf}[BoldFont=lmroman10-bold.otf,
 ItalicFont=lmroman10-italic.otf,BoldItalicFont=lmroman10-bolditalic.otf]
\usepackage{amsmath,amssymb,mathtools}
\usepackage{unicode-math}
\setmathfont{latinmodern-math.otf}
\AtBeginDocument{\let\setminus\smallsetminus}
\usepackage{xurl,hyperref}
\hypersetup{colorlinks=true,urlcolor=blue}
\setcounter{secnumdepth}{0}
\setlength{\parindent}{0pt}
\setlength{\parskip}{5pt}
\setlength{\emergencystretch}{3em}
\begin{document}
\section*{Classify sofic shifts}\label{4600050}
{\small UnsolvedMath ID 4600050 · AMR-045-0050 · Dynamical Systems · AMR Open Problem Lists}
\subsection{Definitions and mathematical statement}
For a finite labelled directed graph $G$, let $X_G$ be the
set of label sequences of all bi-infinite directed paths.
With the product topology and the left shift, this is a
two-sided sofic shift. Equivalently, a sofic shift is a
continuous shift-commuting image of a shift of finite type;
its set of allowed finite words is a regular language.

The target is a complete classification of sofic shifts up
to topological conjugacy. Given finite labelled graph
presentations $G,H$, determine exactly when a homeomorphism
$\Phi:X_G\to X_H$ satisfies
\[
                       \Phi\sigma_G=\sigma_H\Phi.
\]
A useful classification should supply complete, assessable
invariants or necessary and sufficient conditions on the
presentations, including whether conjugacy admits a terminating
decision procedure.

The two coding maps of a conjugacy may use finitely many
past and future coordinates, with no prescribed window size.
They act on label sequences, not on chosen lifts to edge
paths, since distinct paths can have the same labels.
Consequently, equality or isomorphism of the presenting graphs
is not the required equivalence relation. Reducible, periodic,
and mixing sofic systems are all included. The classification
must preserve one application of the shift, rather than only
higher powers, measure-theoretic behavior, or the orbits of
the corresponding suspension flow. Finite-type shifts form
a subclass of this broader target.
\subsection{Short English statement}
Classify two-sided sofic shifts under topological conjugacy, using their finite labelled presentations.
\subsection{Sources}
\begin{itemize}
\item[S1] ULAM.AI and the UnsolvedMath Contributors, supplied problems.json, record 4600050 (AMR-045-0050), AMR Open Problem Lists. Catalogue identity and source transcription; CC BY 4.0.
\url{https://huggingface.co/datasets/ulamai/UnsolvedMath}.
\item[S2] Boyle - Open Problems in Symbolic Dynamics (2008), Problem/Question/Conjecture 32.1. Original source indexed by AMR-045-0050.
\url{https://www.math.umd.edu/~mboyle/papers/openfinalsub3nov2007.pdf}.
\end{itemize}
\end{document}
